% Options for packages loaded elsewhere
\PassOptionsToPackage{unicode}{hyperref}
\PassOptionsToPackage{hyphens}{url}
\PassOptionsToPackage{dvipsnames,svgnames,x11names}{xcolor}
\documentclass[
  11pt,
]{article}
\usepackage{xcolor}
\usepackage[margin=25mm]{geometry}
\usepackage{amsmath,amssymb}
\setcounter{secnumdepth}{-\maxdimen} % remove section numbering
\usepackage{iftex}
\ifPDFTeX
  \usepackage[T1]{fontenc}
  \usepackage[utf8]{inputenc}
  \usepackage{textcomp} % provide euro and other symbols
\else % if luatex or xetex
  \usepackage{unicode-math} % this also loads fontspec
  \defaultfontfeatures{Scale=MatchLowercase}
  \defaultfontfeatures[\rmfamily]{Ligatures=TeX,Scale=1}
\fi
\usepackage{lmodern}
\ifPDFTeX\else
  % xetex/luatex font selection
\fi
% Use upquote if available, for straight quotes in verbatim environments
\IfFileExists{upquote.sty}{\usepackage{upquote}}{}
\IfFileExists{microtype.sty}{% use microtype if available
  \usepackage[]{microtype}
  \UseMicrotypeSet[protrusion]{basicmath} % disable protrusion for tt fonts
}{}
\makeatletter
\@ifundefined{KOMAClassName}{% if non-KOMA class
  \IfFileExists{parskip.sty}{%
    \usepackage{parskip}
  }{% else
    \setlength{\parindent}{0pt}
    \setlength{\parskip}{6pt plus 2pt minus 1pt}}
}{% if KOMA class
  \KOMAoptions{parskip=half}}
\makeatother
\usepackage{listings}
\newcommand{\passthrough}[1]{#1}
\lstset{defaultdialect=[5.3]Lua}
\lstset{defaultdialect=[x86masm]Assembler}
\usepackage{longtable,booktabs,array}
\usepackage{needspace}
\usepackage{calc} % for calculating minipage widths
% Correct order of tables after \paragraph or \subparagraph
\usepackage{etoolbox}
\makeatletter
\patchcmd\longtable{\par}{\if@noskipsec\mbox{}\fi\par}{}{}
\makeatother
% Allow footnotes in longtable head/foot
\IfFileExists{footnotehyper.sty}{\usepackage{footnotehyper}}{\usepackage{footnote}}
\makesavenoteenv{longtable}
\setlength{\emergencystretch}{3em} % prevent overfull lines
\providecommand{\tightlist}{%
  \setlength{\itemsep}{0pt}\setlength{\parskip}{0pt}}
\emergencystretch=3em
\AtBeginDocument{\Urlmuskip=0mu plus 2mu}
\lstset{breaklines=true,basicstyle=\small\ttfamily,columns=fullflexible,keepspaces=true,showstringspaces=false}
\setlength{\tabcolsep}{4pt}
\usepackage{iftex}
\ifXeTeX
  \usepackage{newunicodechar}
  \makeatletter
  \newcommand{\DeclareUnicodeCharacter}[2]{%
    \begingroup
    \uccode`\~="#1\relax
    \uppercase{\def\temp{~}}%
    \expandafter\endgroup
    \expandafter\newunicodechar\expandafter{\temp}{#2}%
  }
  \makeatother
\fi
\ifdefined\DeclareUnicodeCharacter
\DeclareUnicodeCharacter{2192}{\ensuremath{\to}}
\DeclareUnicodeCharacter{2190}{\ensuremath{\leftarrow}}
\DeclareUnicodeCharacter{2194}{\ensuremath{\leftrightarrow}}
\DeclareUnicodeCharacter{21D2}{\ensuremath{\Rightarrow}}
\DeclareUnicodeCharacter{21A6}{\ensuremath{\mapsto}}
\DeclareUnicodeCharacter{2212}{\ensuremath{-}}
\DeclareUnicodeCharacter{2264}{\ensuremath{\leq}}
\DeclareUnicodeCharacter{2265}{\ensuremath{\geq}}
\DeclareUnicodeCharacter{2260}{\ensuremath{\neq}}
\DeclareUnicodeCharacter{2113}{\ensuremath{\ell}}
\DeclareUnicodeCharacter{03A3}{\ensuremath{\Sigma}}
\DeclareUnicodeCharacter{03A6}{\ensuremath{\Phi}}
\DeclareUnicodeCharacter{0393}{\ensuremath{\Gamma}}
\DeclareUnicodeCharacter{221E}{\ensuremath{\infty}}
\DeclareUnicodeCharacter{2208}{\ensuremath{\in}}
\DeclareUnicodeCharacter{2209}{\ensuremath{\notin}}
\DeclareUnicodeCharacter{211A}{\ensuremath{\mathbb{Q}}}
\DeclareUnicodeCharacter{211D}{\ensuremath{\mathbb{R}}}
\DeclareUnicodeCharacter{2124}{\ensuremath{\mathbb{Z}}}
\DeclareUnicodeCharacter{2115}{\ensuremath{\mathbb{N}}}
\DeclareUnicodeCharacter{2200}{\ensuremath{\forall}}
\DeclareUnicodeCharacter{2203}{\ensuremath{\exists}}
\DeclareUnicodeCharacter{2286}{\ensuremath{\subseteq}}
\DeclareUnicodeCharacter{00D7}{\ensuremath{\times}}
\DeclareUnicodeCharacter{03B4}{\ensuremath{\delta}}
\DeclareUnicodeCharacter{2202}{\ensuremath{\partial}}
\DeclareUnicodeCharacter{03B2}{\ensuremath{\beta}}
\DeclareUnicodeCharacter{03B1}{\ensuremath{\alpha}}
\DeclareUnicodeCharacter{03C6}{\ensuremath{\phi}}
\DeclareUnicodeCharacter{03B3}{\ensuremath{\gamma}}
\DeclareUnicodeCharacter{00B9}{\ensuremath{^1}}
\DeclareUnicodeCharacter{00B2}{\ensuremath{^2}}
\DeclareUnicodeCharacter{2070}{\ensuremath{^0}}
\DeclareUnicodeCharacter{2080}{\ensuremath{_0}}
\DeclareUnicodeCharacter{2081}{\ensuremath{_1}}
\DeclareUnicodeCharacter{2082}{\ensuremath{_2}}
\fi
\usepackage{bookmark}
\IfFileExists{xurl.sty}{\usepackage{xurl}}{} % add URL line breaks if available
\urlstyle{same}
\hypersetup{
  pdftitle={Designing executable lessons with explicit evidence boundaries},
  pdfauthor={Jeremy H. Carroll},
  colorlinks=true,
  linkcolor={Maroon},
  filecolor={Maroon},
  citecolor={Blue},
  urlcolor={Blue},
  pdfcreator={LaTeX via pandoc}}

\title{Designing executable lessons with explicit evidence boundaries}
\author{Jeremy H. Carroll}
\date{2026-10-02 · Version 0.1.0 · Curriculum-design preprint · First-publication review}

\AtBeginEnvironment{longtable}{\small}
\begin{document}
\makeatletter
\begin{center}
{\LARGE \@title\par}
\vspace{0.8em}
{\normalsize \@author\par}
\vspace{0.5em}
{\small \@date\par}
\end{center}
\makeatother

Source audit baseline:
\texttt{\nolinkurl{52cb5457f6c485ba548cfd292f6012a941ed22cf}}.
Package reviewed from:
\texttt{\nolinkurl{b31360060128d992113b602adee326e70d223340}}.
External source links resolve at the audit baseline; package-local files
are the revised v0.1.0 candidate.

\hypertarget{abstract}{%
\subsection{Abstract}\label{abstract}}

An executable lesson can run successfully while teaching the wrong
claim. A cell may compare an answer with itself, a notebook may import
an obsolete implementation, or a successful finite check may be
described as general proof. This guide develops a practical lesson
design around five components: a declared finite example, a worked
calculation, a learner prediction, a changed-input probe, and an
explicit evidence boundary. It maps the 20-notebook Symbol Game
integration route at a declared source baseline and works through one
triangle in detail. Source inspection supports the design and its
available artifacts; historical execution records are identified
separately. The package adds a companion notebook that repeats the
worked triangle with exact fractions only, and a pinned Lean project
containing the two formal files the guide cites. No learner study,
attention half-life, teaching-efficiency estimate, or transfer-effect
result is reported.

\hypertarget{what-a-lesson-should-let-the-reader-do}{%
\subsection{1. What a lesson should let the reader
do}\label{what-a-lesson-should-let-the-reader-do}}

The learner's first task is concrete: explain why an outcome changes
when one declared input changes. A \textbf{declared input} is a value or
rule chosen for the example. A \textbf{probe} is a particular changed
input used to test a predicted behavior. A \textbf{reference check} is
the computation against which the prediction is compared. An
\textbf{evidence boundary} states the conclusions that this check
actually supports.

For a first encounter, assume familiarity with integer arithmetic and
introduce fractions before using them. Show what each symbol denotes in
the current example. ``Cochain,'' for instance, can initially mean a
table assigning a rational number to each directed edge. The formal name
becomes useful after the reader has worked with the table.

A complete lesson should let a reader reconstruct the original case,
identify the changed input, predict an outcome with a reason, inspect
the checker, and state one conclusion that the example leaves open. The
notebook can record the reference computation. A record of the learner's
understanding requires the learner's own response and an explicit
assessment procedure.

This distinction addresses concrete defects documented by the existing
\href{https://github.com/dbccpage/the_symbol_game/blob/52cb5457f6c485ba548cfd292f6012a941ed22cf/omega_engine_education/the_symbol_game/lessons/system_integration/curriculum_125_review.md}{curriculum
review}: reused generic arithmetic, constant assertions,
self-comparisons, and matching LP objectives without feasibility checks.
The review describes detected patterns; it does not establish that every
older lesson is wrong.

\hypertarget{a-repeatable-design-sequence}{%
\subsection{2. A repeatable design
sequence}\label{a-repeatable-design-sequence}}

Start with the smallest example that exposes the intended distinction.
Write the claim before writing the notebook. ``A potential change
preserves a directed cycle sum'' is sufficiently precise; ``understand
advanced topology'' is not a checkable outcome for one sitting.

Then give a worked computation with each intermediate value visible. A
learner should be able to follow it without first opening a distant
module to learn what a variable means. Link the real implementation
beside the explanation so that the source remains inspectable.

Next remove a step or vary one input. Ask for a prediction before
showing the reference answer. Follow the successful case with a
plausible failure: a missing bracket, a changed digest-bound atom, a
dual witness that is not closed, or a forbidden numerical type. A useful
probe can fail for a specific reason.

Finally, place the evidence boundary immediately after the result. A
parser checks syntax. A digest check detects an inconsistency between
content and its named digest. A finite certificate checker establishes
its declared property for supplied data. Each deserves its own sentence.

The decision to begin with worked examples and gradually remove guidance
is informed by primary studies of algebra examples and transitions from
example study to problem solving:
\href{https://www.tandfonline.com/doi/abs/10.1207/s1532690xci0201_3}{Sweller
and Cooper (1985)} and
\href{https://doi.org/10.1080/00220970209599510}{Renkl, Atkinson, Maier,
and Staley (2002)}. Applying that sequence here is a curriculum-design
choice. Those studies do not evaluate these notebooks, establish a
universal lesson duration, or measure transfer into formal verification.
The present contribution is the inspectable lesson specification and its
failure controls. It is not a new learning-effect estimate or a claim
that worked examples originated in this series.

\hypertarget{worked-lesson-moving-a-discrepancy-without-erasing-it}{%
\subsection{3. Worked lesson: moving a discrepancy without erasing
it}\label{worked-lesson-moving-a-discrepancy-without-erasing-it}}

This lesson uses a rationally normalized version of the series' (1,1,1)
triangle. It checks cycle sums and residual coordinates; no weights or
quotient-minimum solver are used.

Draw vertices a, b, c with edges a to b, b to c, and c to a. Put \(1/3\)
on every edge. Their signed cycle sum is one. A \textbf{potential}
assigns a height to each vertex; the edge difference is the target
height minus the source height.

Choose heights

\[
f_a=0,\qquad f_b=\tfrac13,\qquad f_c=\tfrac23.
\]

The edge differences are

\[
df=(\tfrac13,\tfrac13,-\tfrac23).
\]

Subtract these differences from the original edge table
\(h=(1/3,1/3,1/3)\). The residual is \(h-df=(0,0,1)\). Its location has
changed, but its cycle sum is still one.

Here is an adapted arithmetic core from the
\href{https://github.com/dbccpage/the_symbol_game/blob/52cb5457f6c485ba548cfd292f6012a941ed22cf/omega_engine_education/the_symbol_game/lessons/system_integration/100_novices/04_cycle_repair.ipynb}{cycle-repair
lesson}, using the repository's discrete cochain implementation. It
omits the notebook's retired \texttt{\nolinkurl{l1_norm()}} call;
Section 5 explains why the unchanged full notebook cannot currently
complete:

\Needspace{17\baselineskip}
\begin{lstlisting}[language=Python]
from fractions import Fraction as Q
from omega_engine_symbol_game.domain.ddc import (
    Edge, SimplicialComplex, Cochain0, Cochain1, d0,
)

edges = (Edge("a", "b"), Edge("b", "c"), Edge("c", "a"))
carrier = SimplicialComplex(("a", "b", "c"), edges)
h = Cochain1({edge.key: Q(1, 3) for edge in edges})
f = Cochain0({"a": Q(0), "b": Q(1, 3), "c": Q(2, 3)})
adjustment = d0(carrier, f)
residual = h - adjustment

assert sum(adjustment.values.values(), Q(0)) == 0
assert sum(residual.values.values(), Q(0)) == 1
assert residual.values[("c", "a")] == 1
\end{lstlisting}

The explanation supplies a reason independent of those assertions:
adding the three endpoint differences cancels every vertex height.
Therefore the sum of \(h-df\) equals the sum of \(h\). The calculation
instantiates that identity on a named triangle.

\hypertarget{prediction-before-the-next-cell}{%
\subsubsection{Prediction before the next
cell}\label{prediction-before-the-next-cell}}

Add \(7/5\) to every vertex height. Predict the new edge differences and
residual before reading the reference check. Identify where the added
quantity appears with opposite signs.

\Needspace{5\baselineskip}
\begin{lstlisting}[language=Python]
shifted = Cochain0({vertex: height + Q(7, 5)
                   for vertex, height in f.values.items()})
assert d0(carrier, shifted) == adjustment
\end{lstlisting}

Both endpoints of every edge receive the same shift, so their difference
is unchanged. Now negate the original edge table while keeping the same
potential. Predict its new signed residual sum.

\Needspace{3\baselineskip}
\begin{lstlisting}[language=Python]
assert sum((-h - adjustment).values.values(), Q(0)) == -1
\end{lstlisting}

The answer changes because the edge claims changed, not because constant
height shifts acquired a new effect. For a further completion exercise,
leave the third original edge value blank and choose it so the cycle sum
becomes zero. For a less guided task, ask the reader to construct
another three-edge example and explain what a potential adjustment can
preserve.

Check: the missing third edge claim is \(-2/3\), since the three claims
must sum to zero. For an independent example, require a stated edge
orientation, exact values, and an explanation that the residual cycle
sum equals the original claim cycle sum. Changing the potential cannot
change that sum.

\hypertarget{a-check-that-can-pass-with-the-wrong-sign}{%
\subsubsection{A check that can pass with the wrong
sign}\label{a-check-that-can-pass-with-the-wrong-sign}}

A cycle-sum assertion alone does not validate the edge convention. A
faulty implementation that returns source height minus target height
also telescopes to zero. On the displayed fixture it gives
\(\widetilde{df}=(-1/3,-1/3,2/3)\), hence
\(h-\widetilde{df}=(2/3,2/3,-1/3)\). This incorrect residual still has
cycle sum one. The predicted coordinates \(df=(1/3,1/3,-2/3)\)
distinguish these implementations. The companion notebook includes this
deliberate wrong-sign control; it requires both the weak check to pass
and the coordinate check to reject the mutant. A control demonstrates
sensitivity to this particular fault, not to every possible
implementation error.

For a learner record, retain the first predicted coordinates and
explanation before revealing outputs, then record any correction
separately. Stored notebook answers and an execution count cannot
reconstruct that first response.

\hypertarget{evidence-boundary}{%
\subsubsection{Evidence boundary}\label{evidence-boundary}}

The computation checks a declared rational triangle and these particular
probes. The telescoping identity explains the cycle invariant. Neither
establishes a physical interpretation, neural behavior, or a learner's
mastery. The code displays residual arithmetic; it does not issue a
canonical native quotient-minimum receipt. Compare the
\href{https://github.com/dbccpage/the_symbol_game/blob/52cb5457f6c485ba548cfd292f6012a941ed22cf/omega_engine_atlas/coda_atlas/omega_engine_symbol_game/005_pfp_practitioners_field_manual/paper.md}{practitioner
guide}, where replacing the claims is a separate operation.

\hypertarget{standalone-replay}{%
\subsubsection{Standalone replay}\label{standalone-replay}}

The package's
\href{jupyter_notebooks/007_moving_a_discrepancy.ipynb}{companion
notebook} repeats this lesson with Python's
\passthrough{\lstinline!fractions!} module only and imports no
repository code. Its revised stored outputs come from a Python 3.12.14
execution.
\texttt{\nolinkurl{tools/replay_companion_notebook.py}} executes
its ordinary Python cells in a fresh CPython process, then compares
sequential code-cell counters and normalized output data with the
deposited ones. It leaves the deposited notebook unchanged and reports
its hash. The stored counts are sequential replay counters, not Jupyter
kernel history. Output equality supports reproducibility of these cells
in that environment; it does not validate their instructional claims. An
optional Jupyter backend requires an explicitly selected kernel. Jupyter
kernel replay could not complete on this review host because sockets are
unavailable; no widget or IPython rendering check is claimed. The
package test
\texttt{\nolinkurl{code/tests/test_007_worked_triangle.py}}
checks the same arithmetic with plain fractions and reads the notebook's
stored outputs without executing it. The same test file also runs the
three code blocks of this section exactly as printed; that check needs
the repository implementation and its initialization dependencies. In a
standalone deposit it is skipped with that reason. The first-publication
review separately executes the printed blocks against seven Git-verified
DDC modules with parent-package initialization isolated; that is a
focused source check, not a full repository integration run.

\hypertarget{the-route-at-the-source-audit-baseline}{%
\subsection{4. The route at the source audit
baseline}\label{the-route-at-the-source-audit-baseline}}

The
\href{https://github.com/dbccpage/the_symbol_game/blob/52cb5457f6c485ba548cfd292f6012a941ed22cf/omega_engine_education/the_symbol_game/lessons/system_integration/README.md}{integration
README} defines the deep route at the source audit baseline. Its five
directory labels organize progression; they are not measured age
thresholds. Begin with the first two groups and use later groups when
their prerequisites are available.

All 20 paths below were verified in the Git trees at the source audit
baseline under
\texttt{\nolinkurl{omega_engine_education/the_symbol_game/lessons/system_integration/}}.
The linked table is a reading map, not a claim of fresh execution.
\texttt{\nolinkurl{docs/source_inventory_2026_10_02.json}}
retains their Git blob identities for inspection after extraction.

\clearpage
\begin{longtable}[]{@{}
  >{\raggedright\arraybackslash}p{(\columnwidth - 4\tabcolsep) * \real{0.1500}}
  >{\raggedright\arraybackslash}p{(\columnwidth - 4\tabcolsep) * \real{0.3800}}
  >{\raggedright\arraybackslash}p{(\columnwidth - 4\tabcolsep) * \real{0.4700}}@{}}
\toprule\noalign{}
\begin{minipage}[b]{\linewidth}\raggedright
Group
\end{minipage} & \begin{minipage}[b]{\linewidth}\raggedright
Notebook
\end{minipage} & \begin{minipage}[b]{\linewidth}\raggedright
Main distinction
\end{minipage} \\
\midrule\noalign{}
\endhead
\bottomrule\noalign{}
\endlastfoot
Beginners &
\href{https://github.com/dbccpage/the_symbol_game/blob/52cb5457f6c485ba548cfd292f6012a941ed22cf/omega_engine_education/the_symbol_game/lessons/system_integration/000_beginners/01_marks_and_packets.ipynb}{\texttt{01\_\allowbreak{}marks\_\allowbreak{}and\_\allowbreak{}packets.ipynb}}
& Parseable syntax versus assigned meaning \\
Beginners &
\href{https://github.com/dbccpage/the_symbol_game/blob/52cb5457f6c485ba548cfd292f6012a941ed22cf/omega_engine_education/the_symbol_game/lessons/system_integration/000_beginners/02_lifts_and_non_descent.ipynb}{\texttt{02\_\allowbreak{}lifts\_\allowbreak{}and\_\allowbreak{}non\_\allowbreak{}descent.ipynb}}
& Preserved old object versus a new object without an old
representative \\
Beginners &
\href{https://github.com/dbccpage/the_symbol_game/blob/52cb5457f6c485ba548cfd292f6012a941ed22cf/omega_engine_education/the_symbol_game/lessons/system_integration/000_beginners/03_edges_and_parallel_cells.ipynb}{\texttt{03\_\allowbreak{}edges\_\allowbreak{}and\_\allowbreak{}parallel\_\allowbreak{}cells.ipynb}}
& Matching boundaries versus valid syntax alone \\
Beginners &
\href{https://github.com/dbccpage/the_symbol_game/blob/52cb5457f6c485ba548cfd292f6012a941ed22cf/omega_engine_education/the_symbol_game/lessons/system_integration/000_beginners/04_costs_and_ties.ipynb}{\texttt{04\_\allowbreak{}costs\_\allowbreak{}and\_\allowbreak{}ties.ipynb}}
& Declared cost choice versus an unavoidable winner \\
Novices &
\href{https://github.com/dbccpage/the_symbol_game/blob/52cb5457f6c485ba548cfd292f6012a941ed22cf/omega_engine_education/the_symbol_game/lessons/system_integration/100_novices/01_codecs_and_tamper_evidence.ipynb}{\texttt{01\_\allowbreak{}codecs\_\allowbreak{}and\_\allowbreak{}tamper\_\allowbreak{}evidence.ipynb}}
& Representation preservation versus changed content \\
Novices &
\href{https://github.com/dbccpage/the_symbol_game/blob/52cb5457f6c485ba548cfd292f6012a941ed22cf/omega_engine_education/the_symbol_game/lessons/system_integration/100_novices/02_token_frontier.ipynb}{\texttt{02\_\allowbreak{}token\_\allowbreak{}frontier.ipynb}}
& Retained token mass versus answer correctness \\
Novices &
\href{https://github.com/dbccpage/the_symbol_game/blob/52cb5457f6c485ba548cfd292f6012a941ed22cf/omega_engine_education/the_symbol_game/lessons/system_integration/100_novices/03_typed_agent_composition.ipynb}{\texttt{03\_\allowbreak{}typed\_\allowbreak{}agent\_\allowbreak{}composition.ipynb}}
& Compatible types versus permission to act \\
Novices &
\href{https://github.com/dbccpage/the_symbol_game/blob/52cb5457f6c485ba548cfd292f6012a941ed22cf/omega_engine_education/the_symbol_game/lessons/system_integration/100_novices/04_cycle_repair.ipynb}{\texttt{04\_\allowbreak{}cycle\_\allowbreak{}repair.ipynb}}
& Residual concentration versus changed cycle sum \\
Intermediates &
\href{https://github.com/dbccpage/the_symbol_game/blob/52cb5457f6c485ba548cfd292f6012a941ed22cf/omega_engine_education/the_symbol_game/lessons/system_integration/200_intermediates/01_nilpotence_and_missing_boundaries.ipynb}{\texttt{01\_\allowbreak{}nilpotence\_\allowbreak{}and\_\allowbreak{}missing\_\allowbreak{}boundaries.ipynb}}
& Complete cell boundaries versus the chain-law calculation \\
Intermediates &
\href{https://github.com/dbccpage/the_symbol_game/blob/52cb5457f6c485ba548cfd292f6012a941ed22cf/omega_engine_education/the_symbol_game/lessons/system_integration/200_intermediates/02_exact_lp_certificates.ipynb}{\texttt{02\_\allowbreak{}exact\_\allowbreak{}lp\_\allowbreak{}certificates.ipynb}}
& Feasible matching witnesses versus equal objective labels \\
Intermediates &
\href{https://github.com/dbccpage/the_symbol_game/blob/52cb5457f6c485ba548cfd292f6012a941ed22cf/omega_engine_education/the_symbol_game/lessons/system_integration/200_intermediates/03_quantifier_order.ipynb}{\texttt{03\_\allowbreak{}quantifier\_\allowbreak{}order.ipynb}}
& Choosing before versus after a challenge \\
Intermediates &
\href{https://github.com/dbccpage/the_symbol_game/blob/52cb5457f6c485ba548cfd292f6012a941ed22cf/omega_engine_education/the_symbol_game/lessons/system_integration/200_intermediates/04_frustrated_spin_triangle.ipynb}{\texttt{04\_\allowbreak{}frustrated\_\allowbreak{}spin\_\allowbreak{}triangle.ipynb}}
& Finite parity obstruction versus a general physics claim \\
Advanced &
\href{https://github.com/dbccpage/the_symbol_game/blob/52cb5457f6c485ba548cfd292f6012a941ed22cf/omega_engine_education/the_symbol_game/lessons/system_integration/300_advanced/01_four_simplex_volume.ipynb}{\texttt{01\_\allowbreak{}four\_\allowbreak{}simplex\_\allowbreak{}volume.ipynb}}
& Computed determinant versus a supplied expected constant \\
Advanced &
\href{https://github.com/dbccpage/the_symbol_game/blob/52cb5457f6c485ba548cfd292f6012a941ed22cf/omega_engine_education/the_symbol_game/lessons/system_integration/300_advanced/02_pairing_and_degree_four_certificate.ipynb}{\texttt{02\_\allowbreak{}pairing\_\allowbreak{}and\_\allowbreak{}degree\_\allowbreak{}four\_\allowbreak{}certificate.ipynb}}
& Declared pairing versus constructed manifold topology \\
Advanced &
\href{https://github.com/dbccpage/the_symbol_game/blob/52cb5457f6c485ba548cfd292f6012a941ed22cf/omega_engine_education/the_symbol_game/lessons/system_integration/300_advanced/03_survivor_birth_and_death.ipynb}{\texttt{03\_\allowbreak{}survivor\_\allowbreak{}birth\_\allowbreak{}and\_\allowbreak{}death.ipynb}}
& Survivor identities versus counts alone \\
Advanced &
\href{https://github.com/dbccpage/the_symbol_game/blob/52cb5457f6c485ba548cfd292f6012a941ed22cf/omega_engine_education/the_symbol_game/lessons/system_integration/300_advanced/04_transfer_trace_and_shadow.ipynb}{\texttt{04\_\allowbreak{}transfer\_\allowbreak{}trace\_\allowbreak{}and\_\allowbreak{}shadow.ipynb}}
& Finite return counts versus broader spectral meaning \\
Researchers &
\href{https://github.com/dbccpage/the_symbol_game/blob/52cb5457f6c485ba548cfd292f6012a941ed22cf/omega_engine_education/the_symbol_game/lessons/system_integration/400_researchers/01_coverage_and_licensed_depth.ipynb}{\texttt{01\_\allowbreak{}coverage\_\allowbreak{}and\_\allowbreak{}licensed\_\allowbreak{}depth.ipynb}}
& Catalog coverage versus accepted lift evidence \\
Researchers &
\href{https://github.com/dbccpage/the_symbol_game/blob/52cb5457f6c485ba548cfd292f6012a941ed22cf/omega_engine_education/the_symbol_game/lessons/system_integration/400_researchers/02_immune_boundary_and_negative_selection.ipynb}{\texttt{02\_\allowbreak{}immune\_\allowbreak{}boundary\_\allowbreak{}and\_\allowbreak{}negative\_\allowbreak{}selection.ipynb}}
& Bounded detection versus repair authority \\
Researchers &
\href{https://github.com/dbccpage/the_symbol_game/blob/52cb5457f6c485ba548cfd292f6012a941ed22cf/omega_engine_education/the_symbol_game/lessons/system_integration/400_researchers/03_governed_prompt_evolution.ipynb}{\texttt{03\_\allowbreak{}governed\_\allowbreak{}prompt\_\allowbreak{}evolution.ipynb}}
& Selection on supplied finite outcomes versus learned model
performance \\
Researchers &
\href{https://github.com/dbccpage/the_symbol_game/blob/52cb5457f6c485ba548cfd292f6012a941ed22cf/omega_engine_education/the_symbol_game/lessons/system_integration/400_researchers/04_source_bound_research_audit.ipynb}{\texttt{04\_\allowbreak{}source\_\allowbreak{}bound\_\allowbreak{}research\_\allowbreak{}audit.ipynb}}
& Source navigation versus executed evidence \\
\end{longtable}

The retained
\href{https://github.com/dbccpage/the_symbol_game/blob/52cb5457f6c485ba548cfd292f6012a941ed22cf/omega_engine_education/the_symbol_game/lessons/curriculum_125_master.md}{125-lesson
syllabus} supplies a broader topic ordering. The integration README
explicitly avoids treating it as a second front door or treating 20
examples as implementation of all 125 topics.

\hypertarget{replay-honestly-after-a-source-move}{%
\subsection{5. Replay honestly after a source
move}\label{replay-honestly-after-a-source-move}}

Execution evidence belongs to exact notebook and implementation bytes.
The saved
\href{https://github.com/dbccpage/the_symbol_game/blob/52cb5457f6c485ba548cfd292f6012a941ed22cf/omega_engine_education/the_symbol_game/lessons/system_integration/headless_receipt.json}{headless
receipt} reports 20 passing notebooks dated 2026-09-25. Its recorded
paths use the former
\texttt{\nolinkurl{omega_engine_education/notebooks/}} tree. The
\href{https://github.com/dbccpage/the_symbol_game/blob/52cb5457f6c485ba548cfd292f6012a941ed22cf/omega_engine_education/the_symbol_game/lessons/system_integration/lesson_manifest.json}{manifest}
also retains 20 paths under that absent location at this baseline. These
are historical records, not a current 20/20 claim.

The
\href{https://github.com/dbccpage/the_symbol_game/blob/52cb5457f6c485ba548cfd292f6012a941ed22cf/omega_engine_education/scripts/verify_system_integration.py}{verifier}
executes ordinary Python cells, with a fresh CPython process per
notebook in its aggregate mode. It does not run an IPython kernel or
test rendered widgets. It checks notebook structure, selected
exact-arithmetic restrictions, source hashes, and runtime assertions. A
cell-count rule cannot independently establish that every cell has the
instructional role its heading claims.

The cycle notebook has an additional, concrete blocker: it calls
\texttt{\nolinkurl{residual.l1_norm()}}, while the pinned
\href{https://github.com/dbccpage/the_symbol_game/blob/52cb5457f6c485ba548cfd292f6012a941ed22cf/omega_engine_symbol_game/domain/ddc/cochains.py}{cochain
implementation} unconditionally raises
\texttt{\nolinkurl{FunctionalBoundaryUnavailable}} for that method.
An explicit-path diagnostic bypasses the stale aggregate manifest but is
expected to stop at this retired API if earlier source checks succeed:

\Needspace{7\baselineskip}
\begin{lstlisting}
python -m pip install -r omega_engine_education/the_symbol_game/lessons/system_integration/requirements.txt
python -B omega_engine_education/scripts/verify_system_integration.py --child omega_engine_education/the_symbol_game/lessons/system_integration/100_novices/04_cycle_repair.ipynb --output "$env:TEMP/pfp-cycle-lesson.json"
\end{lstlisting}

Use Python 3.12 or newer as the route specifies. The child command
validates declared source paths and writes its selected receipt to the
requested output path. Report the observed failure; do not relabel it a
pass because the mathematics is sound. Section 3's adapted arithmetic
core is the available example to check separately. No fresh execution of
this diagnostic or the complete route is reported by the
first-publication review.

A full route refresh needs reconciled manifest/source paths, maintained
replacements for retired API calls, and a new receipt. Preserve the
older receipt's identity. The deposited companion uses its own
maintained arithmetic lesson; refreshing the external route remains
separate work.

\hypertarget{audit-the-lessons-strongest-sentence}{%
\subsection{6. Audit the lesson's strongest
sentence}\label{audit-the-lessons-strongest-sentence}}

Ask what would make the lesson's headline false. A codec lesson should
alter content while keeping the old digest; checking only digest length
misses the intended failure. An LP lesson should reject an infeasible
witness even when the displayed objective matches. A geometry lesson
should compute the determinant independently, rather than assign the
expected determinant and compare it with itself.

Formal references need the same precision.
\href{lean4/GrammarCeiling.lean}{GrammarCeiling.lean} formalizes a bound
for an integer list with both a length bound and a per-entry magnitude
bound. It does not establish an unrestricted cognitive workload.
\href{lean4/GrammarDepthCeiling.lean}{GrammarDepthCeiling.lean} shows
that repeated unary nesting preserves supplied declaration predicates;
it does not bound nesting depth or student search.

This package vendors both files, with the one module the first imports
(\href{lean4/Expressivity.lean}{Expressivity.lean}), into a local Lean
project in \texttt{\nolinkurl{lean4/}} pinned to Lean 4.34.0. Each
vendored file differs from its source at the audited commit only by a
leading provenance comment. The project uses no external library. The
historical \href{lean4/axiom_report.txt}{named axiom report} covers 21
explicitly selected declarations: 19 list no axioms, and the two list
bounds list \passthrough{\lstinline!propext!} and
\passthrough{\lstinline!Quot.sound!}. It is not a census of every
environment declaration. The historical
\href{lean4/full_axiom_report.txt}{full report} lists 64 package-local
declarations, including generated and private names; its recorded
dependencies are subsets of \passthrough{\lstinline!propext!} and
\passthrough{\lstinline!Quot.sound!}.
\passthrough{\lstinline!Classical.choice!} is permitted by the report's
whitelist but is not listed as used. Imported Lean library declarations
are outside that package-local census.

The retained 2026-10-02 receipt binds the full report and formal source
hashes to a successful build and census run. This review verified those
hashes and the vendored source bodies, and checked generation of the
named census. It did not repeat the Lean build because Lean/Lake are
unavailable on the review host. No fresh kernel-verification claim
follows from reading a saved report. These statements concern the
supplied finite formal models and provide no learner result.

Likewise, the package's
\href{code/tests/test_007_pedagogical_transmission.py}{legacy
pedagogical test module} checks chosen timing/metric thresholds and
counts supplied execution-status fields. Those small examples are not
measurements of attention, evidence of a 90\% learner pass rate, or
execution of the 20-notebook route. The current guide makes none of
those claims.

\hypertarget{what-an-educational-study-would-add}{%
\subsection{7. What an educational study would
add}\label{what-an-educational-study-would-add}}

The
\href{https://github.com/dbccpage/the_symbol_game/blob/52cb5457f6c485ba548cfd292f6012a941ed22cf/omega_engine_education/pedagogy/teaching_and_learning_findings.md}{pedagogy
audit} explicitly reports no learner corpus connecting participants with
notebook assertions. Its synthesized research narrative failed citation
validation, so it is used here as a repository audit pointer rather than
as a validated literature synthesis.

A future study could compare lesson formats on predefined graph and
certificate tasks. It would need participant recruitment and consent,
baseline knowledge, recorded first predictions, consistent feedback,
independent assessment, and a delayed transfer task. Timing, exclusions,
missing responses, and analysis decisions would need specification
before examining results. The comparison should test an instructional
choice rather than reward familiarity with local vocabulary.

Such a study remains proposed. The useful present contribution is a
lesson-authoring method and a navigable set of finite examples with
auditable boundaries. Software replay can check the examples; assessing
what people learned is a separate investigation.

\hypertarget{verification-for-this-revision}{%
\subsection{Verification for this
revision}\label{verification-for-this-revision}}

The first-publication record is
\url{docs/first_publication_validation_2026_10_02.json}. It binds the
tested package and source bytes to the observed commands and separates
standalone tests, focused DDC execution, notebook replay, and historical
Lean evidence. The \href{docs/first_publication_review.md}{review}
records corrections and remaining limits.

\texttt{\nolinkurl{verify_deposit.py}} checks file hashes and a
two-way census by default; it does not run tests or Lean unless
explicitly requested.
\texttt{tools/rebuild\_release.py --check} checks the
archive against both manifests. File identity is an integrity check, not
proof of the manuscript's conclusions. Historical receipts retain their
original dates, commands, and scopes.

The \href{_lev_landau_theoretical_minimum_quiz/study_guide.md}{Landau
study guide} provides worked examples, a standard-library worksheet, and
a prediction record. The
\href{_lev_landau_theoretical_minimum_quiz/quiz.md}{24-problem quiz}
progresses from the triangle to checker design, source custody, formal
assumptions, and a proposed learner study; the
\href{_lev_landau_theoretical_minimum_quiz/instructor_key.md}{instructor
key} gives explicit scoring. Its levels and readiness scores are
instructional choices, not validated learner classifications.

\hypertarget{references}{%
\subsection{References}\label{references}}

\begin{itemize}
\tightlist
\item
  Sweller, J., and Cooper, G. A. (1985).
  \href{https://doi.org/10.1207/s1532690xci0201_3}{The Use of Worked
  Examples as a Substitute for Problem Solving in Learning Algebra}.
  \emph{Cognition and Instruction}, 2(1), 59--89. DOI
  10.1207/s1532690xci0201\_3.
\item
  Renkl, A., Atkinson, R. K., Maier, U. H., and Staley, R. (2002).
  \href{https://doi.org/10.1080/00220970209599510}{From Example Study to
  Problem Solving: Smooth Transitions Help Learning}. \emph{The Journal
  of Experimental Education}, 70(4), 293--315. DOI
  10.1080/00220970209599510.
\end{itemize}

\hypertarget{companion-use}{%
\subsection{Companion use}\label{companion-use}}

\begin{itemize}
\tightlist
\item
  \href{https://github.com/dbccpage/the_symbol_game/blob/52cb5457f6c485ba548cfd292f6012a941ed22cf/omega_engine_atlas/coda_atlas/omega_engine_symbol_game/005_pfp_practitioners_field_manual/paper.md}{Companion
  005} supplies the operator route.
\item
  \href{https://github.com/dbccpage/the_symbol_game/blob/52cb5457f6c485ba548cfd292f6012a941ed22cf/omega_engine_atlas/coda_atlas/omega_engine_symbol_game/009_pfp_adversarial_dialectics_and_terminal_reality_gates/paper.md}{Companion
  009} supplies the shared claim-review protocol.
\end{itemize}

\end{document}
